\documentclass[11pt]{amsart}
\usepackage{graphicx}
\usepackage{subfigure}
\usepackage{amssymb}
\usepackage{multicol}
\usepackage{fancyhdr}
\usepackage{answers}
\usepackage{verbatim}

\textwidth = 6.5 in
\textheight = 8.5 in
\oddsidemargin = 0.0 in
\evensidemargin = 0.0 in
\topmargin = 0.0 in
\headheight = 0.0 in
\headsep = 0.2 in
\footskip = 0.3 in
\parskip = 0.2in
\parindent = 0.0in

\newtheorem{theorem}{Theorem}
\newtheorem{corollary}[theorem]{Corollary}
\newtheorem{definition}{Definition}

\newenvironment{packed_enum}{
\begin{enumerate}
  \setlength{\itemsep}{1pt}
  \setlength{\parskip}{0pt}
  \setlength{\parsep}{0pt}
}{\end{enumerate}}

\newenvironment{packed_item}{
\begin{itemize}
  \setlength{\itemsep}{1pt}
  \setlength{\parskip}{0pt}
  \setlength{\parsep}{0pt}
}{\end{itemize}}


\newcounter{homeworkcounter}
\newcounter{problemcounter}
%\newcounter{proofpcounter}

\newenvironment{homework}%
{\refstepcounter{homeworkcounter}
In order to receive full credit problems and solutions must be clearly written and presented neatly.  In writing up your solution, you should state the problem as explaining the solution.  Merely giving answers to exercises will is not enough to receive credit.  Your write up must show work and convince me that you understand the material.  Late homework will not be graded.  

\begin{packed_item}
\item Choose, complete, and write-up 6 problems each assignment.  
\item Put your name on each page you submit, and somewhere indicate who you worked with on this assignment.
\end{packed_item}
}{\clearpage \setcounter{problemcounter}{0} }


\newenvironment{problem}%
{\refstepcounter{problemcounter}\vspace{10pt}\par\noindent
\textbf{Problem \thehomeworkcounter.\theproblemcounter}: }%
{\vspace{0cm}}

%\newenvironment{proofp}%
%{\refstepcounter{proofpcounter}\vspace{10pt}\par\noindent
%\textbf{Proof \thehomeworkcounter.\theproofpcounter}: }%
%{\vspace{0cm}}




\begin{document}
\setcounter{homeworkcounter}{0}

\thispagestyle{fancy}
\pagestyle{fancy}
\fancyhead[CE,CO]{ {\bf Homework \thehomeworkcounter} - Math 393 - Abstract Algebra I}
\fancyfoot[C]{\thepage}



%Homework 1
\begin{homework}
This homework is based on material from Chapter 3 of the text, specifically examples and the definition of groups.
\begin{problem}
Here is a placeholder for a problem.
\end{problem}


\end{homework}

%Homework 2
\begin{homework}
This homework is based on material from Chapter 3 of the text.
\begin{problem}
Here is a placeholder for a problem.
\end{problem}


\end{homework}

%Homework 3
\begin{homework}
This homework is based on material from Chapter 4 of the text.
\begin{problem}
Here is a placeholder for a problem.
\end{problem}


\end{homework}

%Homework 4
\begin{homework}
This homework is based on material from Chapter 5 of the text.
\begin{problem}
Here is a placeholder for a problem.
\end{problem}


\end{homework}

%Homework 5
\begin{homework}
This homework is based on material from Chapter 6 of the text, specifically examples and the definition of cosets.
\begin{problem}
Here is a placeholder for a problem.
\end{problem}


\end{homework}

%Homework 6
\begin{homework}
This homework is based on material from Chapter 6 of the text.
\begin{problem}
Here is a placeholder for a problem.
\end{problem}


\end{homework}

%Homework 7
\begin{homework}
This homework is based on material from Chapter 9.
\begin{problem}
Here is a placeholder for a problem.
\end{problem}


\end{homework}

%Homework 8
\begin{homework}
This homework is based on material from Chapter 10 of the text, specifically Factor Group and Normal Groups.
\begin{problem}
Here is a placeholder for a problem.
\end{problem}


\end{homework}

%Homework 9
\begin{homework}
This homework is based on material from Chapter 10 of the text.
\begin{problem}
Here is a placeholder for a problem.
\end{problem}


\end{homework}

%Homework 10
\begin{homework}
This homework is based on material from Chapter 11 of the text.
\begin{problem}
Here is a placeholder for a problem.
\end{problem}


\end{homework}

%Homework 11
\begin{homework}
This homework is based on material from Chapter 16 of the text, specifically examples and the definition of rings, integral domains, and fields.
\begin{problem}
Here is a placeholder for a problem.
\end{problem}


\end{homework}

%Homework 12
\begin{homework}
This homework is based on material from Chapter 16 of the text, specifically ring homomorphisms.
\begin{problem}
Here is a placeholder for a problem.
\end{problem}


\end{homework}

%Homework 13
\begin{homework}
This homework is based on material from Chapter 16 of the text.
\begin{problem}
Here is a placeholder for a problem.
\end{problem}


\end{homework}

\end{document}